% NOTE: Switch to 'acmart' when texlive-publishers is available:
%   \documentclass[acmsmall,screen,review]{acmart}
\documentclass[11pt,a4paper]{article}
\usepackage[margin=1in]{geometry}
\usepackage{hyperref}
\usepackage{authblk}
% Stubs for acmart commands not in article class
\newcommand{\affiliation}[1]{}
\newcommand{\subtitle}[1]{}
\newcommand{\keywords}[1]{\par\smallskip\noindent\textbf{Keywords:} #1}
\newcommand{\email}[1]{}

% --- Packages ---
\usepackage{amsmath,amssymb,amsthm}
\usepackage{mathtools}
\usepackage{stmaryrd}
\usepackage{tikz}
\usepackage{listings}
\usepackage{xcolor}
\usepackage{booktabs}
\usepackage{subcaption}

% --- Theorem environments ---
\newtheorem{theorem}{Theorem}[section]
\newtheorem{lemma}[theorem]{Lemma}
\newtheorem{corollary}[theorem]{Corollary}
\newtheorem{proposition}[theorem]{Proposition}
\theoremstyle{definition}
\newtheorem{definition}[theorem]{Definition}
\newtheorem{example}[theorem]{Example}
\theoremstyle{remark}
\newtheorem{remark}[theorem]{Remark}

% --- Notation ---
\newcommand{\NN}{\mathbb{N}}
\newcommand{\ZZ}{\mathbb{Z}}
\newcommand{\RR}{\mathbb{R}}
\newcommand{\CC}{\mathbb{C}}
\newcommand{\HH}{\mathbb{H}}
\newcommand{\OO}{\mathbb{O}}
\newcommand{\SSS}{\mathbb{S}}
\newcommand{\CD}{\mathrm{CD}}
\newcommand{\Cl}{\mathrm{Cl}}
\newcommand{\PG}{\mathrm{PG}}
\newcommand{\PSL}{\mathrm{PSL}}
\newcommand{\Fano}{\mathcal{F}}
\newcommand{\NonFano}{\overline{\mathcal{F}}}
\newcommand{\synth}{\Rightarrow}
\newcommand{\chk}{\Leftarrow}
\newcommand{\wave}{\mathrm{wave}}
\newcommand{\assocR}{\alpha_R}
\newcommand{\assocL}{\alpha_L}
\newcommand{\dagop}{{\dagger}}
\DeclareMathOperator{\sgn}{sgn}

% --- Code listing style ---
\lstdefinelanguage{Sounio}{
  keywords={fn,let,var,if,else,while,return,assert,with,struct,enum,match},
  keywordstyle=\bfseries\color{blue!70!black},
  commentstyle=\itshape\color{gray},
  stringstyle=\color{red!60!black},
  morecomment=[l]{//},
  sensitive=true,
  basicstyle=\small\ttfamily,
  columns=flexible,
  breaklines=true,
}
\lstdefinelanguage{Lean4}{
  keywords={def,theorem,by,native_decide,where,deriving,inductive,structure,namespace,open,import,let,if,then,else,fun,match,return,with},
  keywordstyle=\bfseries\color{purple!70!black},
  commentstyle=\itshape\color{gray},
  morecomment=[l]{--},
  morecomment=[s]{/-}{-/},
  sensitive=true,
  basicstyle=\small\ttfamily,
  columns=flexible,
  breaklines=true,
}

% --- Metadata ---
\title{Fano-Selective Reassociation: Cocycle Coherence and\\Compiler Rewrite Discipline for Non-Associative Types}
\subtitle{With Formal Verification in Lean~4 and Implementation in Sounio}

\author{Demetrios Agourakis}
\affiliation{%
  \institution{PUC-SP / S\~{a}o Leopoldo Mandic}
  \country{Brazil}
}
\email{demetrios@example.com}

\begin{document}
\maketitle

\begin{abstract}
The Cayley--Dickson octonion algebra has two layers of structure
relevant to compiler optimisation.
\emph{Categorically}, its reassociation law is a cocycle-coherent,
everywhere-invertible associator: the sign factor
$\varphi \in \{+1,-1\}$ satisfies the 3-cocycle condition of
Albuquerque--Majid and thus the pentagon axiom of Lack--Street's
skew monoidal framework.
\emph{Operationally}, only a Fano-selected subset of these
reassociations are semantics-preserving compiler rewrites: exactly
175 of 343 ordered triples of imaginary basis elements yield
identical results under both parenthesisations (42 from Fano-line
collinearity in $\PG(2,2)$ and 133 from repeated-index
alternativity), while the remaining 168 triples change sign, with
$168 = |\PSL(2,7)|$.

We formalise this two-layer structure and implement it in Sounio,
a self-hosted programming language. The type checker requires the
\texttt{NonAssoc} algebraic effect for non-associative
multiplication; the e-graph optimiser uses the Fano predicate to
decide which reassociations are safe. The 168 sign-changing triples
decompose under a dagger involution into $84 + 84$
forward/backward pairs with zero fixed points. All results are
formalised in Lean~4 (95 theorems, no \texttt{sorry}, no Mathlib
dependency) and cross-validated against the running compiler.
\end{abstract}

\keywords{non-associative algebra, algebraic effects, Fano plane,
  skew monoidal category, Cayley--Dickson, formal verification, Lean~4}

% ================================================================
\section{Introduction}\label{sec:intro}
% ================================================================

The associator of an algebra measures the failure of associativity:
\[
  [a, b, c] \;=\; (a \cdot b) \cdot c \;-\; a \cdot (b \cdot c).
\]
In an associative algebra the associator vanishes identically, and a
compiler may freely rewrite $(a \cdot b) \cdot c$ as $a \cdot (b \cdot c)$.
In a non-associative algebra it does not, and the compiler must
preserve the programmer's chosen parenthesisation.

But in the Cayley--Dickson algebras---the tower
$\RR \to \CC \to \HH \to \OO \to \SSS \to \cdots$ that extends the
real numbers through complex numbers, quaternions, and
octonions---the situation is more subtle. Not all triples are
non-associative. For the octonions ($k=3$), exactly 175 of 343
ordered triples of imaginary basis elements are associative, and 168
are not. Among the 210 triples with three distinct indices, the
associative ones are precisely those whose indices form a line of the
\emph{Fano plane} $\PG(2,2)$: $[e_i, e_j, e_l] = 0$ for distinct
nonzero $i,j,l$ iff $\{i,j,l\}$ is one of the 7 lines of $\PG(2,2)$
\cite{baez2002octonions}. The remaining 133 triples (with a repeated
index) are associative by the alternativity of the octonions.

This paper introduces \emph{Fano-selective reassociation}, a
two-layer structure for non-associative type systems. Categorically,
the octonion reassociation is cocycle-coherent and invertible: the
sign factor $\varphi(a,b,c) \in \{+1,-1\}$ satisfies the 3-cocycle
condition of Albuquerque--Majid \cite{albuquerque1999} and thus the
pentagon axiom of the skew monoidal framework
\cite{szlachanyi2012,lackstreet2012} (where associators are not
required to be isomorphisms, though they may be). Operationally,
only the 175 reassociations with $\varphi = +1$ are
semantics-preserving compiler rewrites; the 168 with $\varphi = -1$
change the sign of the computed value. We define a \emph{rewrite
category} whose morphisms are the sign-preserving reassociations,
and show that the Fano plane determines which morphisms exist.

\paragraph{Contributions.}
\begin{enumerate}
\item \textbf{Two-layer structure} (Definition~\ref{def:fsi},
  Theorem~\ref{thm:coherence}): a cocycle-coherent associator
  at the algebraic level, and a Fano-selective rewrite discipline
  at the compiler level, with the Fano plane of $\PG(2,2)$
  determining which rewrites are semantics-preserving
  (\S\ref{sec:fsi}).

\item \textbf{The 168 Theorem, formalised} (\S\ref{sec:formal}):
  95 theorems in Lean~4, including the exact counts $175 + 168 = 343$,
  the binary norm property, and the $84/84$ dagger decomposition.
  No \texttt{sorry}. No Mathlib dependency. Build time under 3 seconds.

\item \textbf{The \texttt{NonAssoc} algebraic effect} (\S\ref{sec:impl}):
  an implementation in Sounio where the type checker requires the
  \texttt{NonAssoc} effect for multiplication in non-associative
  algebras ($k \geq 3$), and the optimiser uses the ternary Fano
  predicate to decide reassociation safety for specific basis triples.
  This makes $\PG(2,2)$ an operational part of the compiler.

\item \textbf{Structural analogy to the two-state vector formalism}
  (\S\ref{sec:tsvf}): we observe that the $84/84$ dagger decomposition
  exhibits the same symmetry as the forward/backward state vectors
  in the Aharonov--Bergmann--Lebowitz time-symmetric, pre- and
  post-selected quantum formalism \cite{abl1964}. We present this as
  a suggestive structural parallel, not a physical claim.
\end{enumerate}

% ================================================================
\section{Background}\label{sec:background}
% ================================================================

\subsection{Cayley--Dickson Algebras and the Sign Function}

The Cayley--Dickson construction \cite{dickson1919} builds a tower of
$2^k$-dimensional algebras:
\[
  \RR \;\xrightarrow{k=1}\; \CC \;\xrightarrow{k=2}\; \HH
  \;\xrightarrow{k=3}\; \OO \;\xrightarrow{k=4}\; \SSS
  \;\xrightarrow{}\; \cdots
\]
At each step, commutativity or associativity is lost. The multiplication
of basis elements is determined by a recursive \emph{sign function}
$\sigma: \{0,\ldots,2^k-1\}^2 \times \NN \to \{+1,-1\}$:
\[
  e_a \cdot e_b \;=\; \sigma(a,b,k) \cdot e_{a \oplus b}
\]
where $\oplus$ denotes bitwise XOR. The sign function satisfies the
recursive Cayley--Dickson doubling formula (see Ren \& Zhao
\cite{renzhao2023} for the explicit twist function).

The \emph{associator} of basis elements $(e_i, e_j, e_l)$ decomposes as:
\begin{align}
  [e_i, e_j, e_l] &= (\alpha - \beta) \cdot e_{i \oplus j \oplus l}
  \label{eq:assoc-decomp}
\end{align}
where $\alpha = \sigma(i,j)\,\sigma(i \oplus j, l)$ and
$\beta = \sigma(j,l)\,\sigma(i, j \oplus l)$. Since
$\alpha, \beta \in \{+1,-1\}$, we have $\alpha - \beta \in \{-2, 0, +2\}$.
We write $\wave(i,j,l) := \alpha - \beta$ for the \emph{associator
coefficient} of the triple.

\begin{theorem}[Two-Valued Associator]\label{thm:binary-norm}
For all $k \geq 3$ and all nonzero $i, j, l \in \{1,\ldots,2^k-1\}$:
the associator coefficient $\alpha - \beta \in \{-2, 0, +2\}$,
so $\lVert[e_i, e_j, e_l]\rVert \in \{0, 2\}$.
\end{theorem}

For $k=3$ (octonions), the $7^3 = 343$ ordered triples of imaginary
basis elements partition as follows:

\begin{center}
\begin{tabular}{lrl}
\toprule
Class & Count & Condition \\
\midrule
Fano-line (distinct, collinear) & 42 & $i,j,l$ distinct, collinear in $\PG(2,2)$ \\
Repeated-index (alternativity) & 133 & $i{=}j$ or $j{=}l$ or $i{=}l$ \\
\cmidrule{2-2}
\emph{Associative subtotal} & 175 & $\alpha = \beta$ \\
\midrule
Forward (non-associative) & 84 & $\alpha - \beta = +2$ \\
Backward (non-associative) & 84 & $\alpha - \beta = -2$ \\
\cmidrule{2-2}
\emph{Non-associative subtotal} & 168 & $\alpha \neq \beta$ \\
\midrule
Total & 343 & \\
\bottomrule
\end{tabular}
\end{center}

The number 168 $= |\PSL(2,7)|$ is the order of the automorphism group
of the Fano plane $\PG(2,2)$ \cite{baez2002octonions}.

\subsection{Skew Monoidal Categories}

A \emph{skew monoidal category} \cite{szlachanyi2012,lackstreet2012}
is a category $\mathcal{C}$ equipped with a bifunctor
$\otimes: \mathcal{C} \times \mathcal{C} \to \mathcal{C}$, a unit
object $I$, and natural transformations
\begin{align*}
  \alpha_{A,B,C} &: (A \otimes B) \otimes C \to A \otimes (B \otimes C) \\
  \lambda_A &: I \otimes A \to A \\
  \rho_A &: A \to A \otimes I
\end{align*}
that are \emph{not} required to be isomorphisms, subject to five
coherence conditions. When all structure maps are invertible, this
recovers an ordinary monoidal category.

The key insight: in a skew monoidal category, re-parenthesisation is a
\emph{directed} operation. The Tamari lattice \cite{tamari1962} on
binary trees makes this precise: right-rotation
$((a \cdot b) \cdot c) \mapsto (a \cdot (b \cdot c))$ is a valid step,
but the reverse is not.

\subsection{Dagger Categories}

A \emph{dagger category} \cite{selinger2007} is a category
$\mathcal{C}$ with an involutive, identity-on-objects, contravariant
functor $\dagop: \mathcal{C}^{\mathrm{op}} \to \mathcal{C}$. For every
morphism $f: A \to B$, there exists $f^\dagop: B \to A$ with
$(f^\dagop)^\dagop = f$.

In categorical quantum mechanics \cite{abramsky2004}, the dagger
models time-reversal: reflecting a string diagram about a horizontal
axis gives the dagger of the process.

\subsection{Algebraic Effects}

Algebraic effects \cite{plotkin2009} provide a modular account of
computational side effects. A function's type signature declares which
effects it may perform:
\[
  f: \tau_1 \to \tau_2 \;\mathbf{with}\; \rho
\]
where $\rho$ is an \emph{effect row}---a set of named effects. Effect
rows form a Boolean algebra under union, intersection, and complement,
with subrow inclusion providing effect subtyping
\cite{leijen2014,lindley2012}.

% ================================================================
\section{Fano-Selective Reassociation: Two Layers}\label{sec:fsi}
% ================================================================

\begin{definition}[Fano-selective reassociation structure]\label{def:fsi}
Let $A$ be a $2^k$-dimensional Cayley--Dickson algebra with basis
$\{e_0, \ldots, e_{2^k-1}\}$, product
$e_i \cdot e_j = \sigma(i,j,k) \cdot e_{i \oplus j}$. Let
$B_0 = \{0, \ldots, 2^k - 1\}$ index all basis elements and
$B = \{1, \ldots, 2^k - 1\}$ the imaginary ones.
A \emph{Fano-selective reassociation structure} on $A$ consists of
two layers:

\paragraph{Algebraic layer (cocycle).}
\begin{enumerate}
\item A reassociation factor
  $\varphi: B_0^3 \to \{+1,-1\}$ defined by
  $\varphi(i,j,l) = \alpha(i,j,l) \cdot \beta(i,j,l)$, where
  $\alpha$ and $\beta$ are the left and right parenthesisation signs.
  This factor is a normalised 3-cocycle on $\ZZ_2^k$
  (\S\ref{sec:coherence}).

\item For every triple $(i,j,l) \in B_0^3$, a signed reassociation
  map
  \[
    \assocR_{i,j,l}: (e_i \cdot e_j) \cdot e_l \;\longmapsto\;
    \varphi(i,j,l) \cdot e_i \cdot (e_j \cdot e_l)
  \]
  that is always invertible as a linear map (since $\varphi^{-1} = \varphi$).
\end{enumerate}

\paragraph{Rewrite layer (compiler).}
\begin{enumerate}\setcounter{enumi}{2}
\item A decidable predicate $\Fano: B^3 \to \{0,1\}$, defined by
  $\Fano(i,j,l) = 1$ iff $\varphi(i,j,l) = +1$.

\item For \emph{associative} triples ($\Fano(i,j,l) = 1$), the two
  parenthesisations yield identical results, so the reassociation is a
  valid semantics-preserving compiler rewrite.

\item For \emph{sign-changing} triples ($\Fano(i,j,l) = 0$), the
  reassociation changes the computed value by a factor of $-1$, so
  no semantics-preserving rewrite exists. These triples have no
  morphism in the \emph{rewrite category} (\S\ref{sec:rewrite-cat}).
\end{enumerate}

A \emph{dagger} on such a structure is an involution
$\dagop: B^3 \to B^3$, $(i,j,l) \mapsto (l,j,i)$, such that:
\begin{itemize}
\item The associative set is dagger-closed:
  $\Fano(i,j,l) = 1 \implies \Fano(l,j,i) = 1$.
\item The dagger negates the associator coefficient on sign-changing
  triples: $\Fano(i,j,l) = 0 \implies \wave(l,j,i) = -\wave(i,j,l)$.
\end{itemize}
\end{definition}

\begin{remark}
The algebraic layer is an everywhere-invertible cocycle associator
in the sense of Albuquerque--Majid \cite{albuquerque1999}. The
rewrite layer extracts from it a selective discipline: the compiler
treats sign-preservation as an effect boundary, blocking
reassociations that change the computed value. These are genuinely
two different categories (\S\ref{sec:coherence},
\S\ref{sec:rewrite-cat}).
\end{remark}

\begin{remark}
When $\Fano \equiv 1$ (all triples), every reassociation is valid and
the algebra is associative. When $\Fano \equiv 0$, no reassociation is
safe. The Fano-selective structure interpolates between these extremes.
The relationship to the partially normal skew monoidal categories
of Uustalu--Veltri--Zeilberger \cite{uvz2021}, where selected
structural laws are invertible, deserves further investigation.
\end{remark}

\subsection{The Octonionic Instance}

For the Cayley--Dickson algebra at tower level $k$, define
$\Fano(i,j,l) = 1$ iff $\alpha(i,j,l) = \beta(i,j,l)$, where
$\alpha$ and $\beta$ are the left and right parenthesisation signs
from \eqref{eq:assoc-decomp}.

\begin{theorem}[The 168 Theorem]\label{thm:168}
For $k = 3$ (octonions):
\begin{enumerate}
\item $|\{(i,j,l) : \Fano(i,j,l) = 1\}| = 175$.
\item $|\{(i,j,l) : \Fano(i,j,l) = 0\}| = 168 = |\PSL(2,7)|$.
\item $175 + 168 = 343 = 7^3$.
\end{enumerate}
\end{theorem}

\begin{proof}
Among the $7^3 = 343$ ordered triples, exactly $343 - 7 \cdot 6 \cdot 5 = 133$
have a repeated index; these are all associative by the alternativity of
the octonions (any subalgebra generated by two elements is associative,
by Artin's theorem). Among the $210$ triples with three distinct indices,
each of the 7 Fano lines contributes $3! = 6$ ordered triples, giving
$42$ associative distinct triples. The remaining $210 - 42 = 168$
distinct triples are non-associative. Hence $175 = 133 + 42$ and
$168 = 210 - 42$. This count coincides with the well-known order
$|\PSL(2,7)| = 168$ of the automorphism group of $\PG(2,2)$.

Independently verified in Lean~4 via \texttt{native\_decide}
(\texttt{fano\_count\_175}, \texttt{non\_fano\_count\_168} in
\texttt{SounioCayleyDickson.lean}).
\end{proof}

\subsection{The Dagger Decomposition}

Define the dagger on triples by reversal: $\dagop(i,j,l) = (l,j,i)$.

\begin{theorem}[$84/84$ Duality]\label{thm:8484}
The 168 non-Fano triples decompose under the dagger into:
\begin{enumerate}
\item 84 \emph{forward} triples with $\wave(i,j,l) = +2$.
\item 84 \emph{backward} triples with $\wave(i,j,l) = -2$.
\item For every forward triple $(i,j,l)$, the dagger
  $(l,j,i)$ is backward: $\wave(l,j,i) = -\wave(i,j,l)$.
\item No non-Fano triple is a fixed point of the dagger.
\end{enumerate}
\end{theorem}

\begin{proof}
By exhaustive computation (theorems \texttt{arrow\_forward\_84},
\texttt{arrow\_backward\_84}, \texttt{dagger\_reversal}, and
\texttt{no\_nonfano\_self\_dual} in \texttt{SounioCayleyDickson.lean}).
\end{proof}

\begin{corollary}\label{cor:tripartition}
The 343 imaginary triples admit a tripartition:
\[
  \underbrace{175}_{\text{invertible}} \;+\;
  \underbrace{84}_{\text{forward}} \;+\;
  \underbrace{84}_{\text{backward}} \;=\; 343
\]
with the dagger preserving the associative set \emph{setwise}
(not pointwise: $(i,j,l)$ and $(l,j,i)$ are generally distinct
ordered triples) and acting as a free involution on the $168$
non-associative triples.
\end{corollary}

\subsection{The Tamari Lattice and Selective Invertibility}

The Tamari lattice $T_n$ \cite{tamari1962} is the partial order on
binary trees with $n$ leaves, where the covering relation is a single
right-rotation: $((a \cdot b) \cdot c) \leq (a \cdot (b \cdot c))$.

In a Fano-selective structure, each edge of the Tamari lattice
corresponding to a Fano triple becomes \emph{undirected} (the rotation
is invertible), while edges corresponding to non-Fano triples remain
\emph{directed}. The result is a partially directed Tamari lattice---a
Hasse diagram where some edges are arrows and others are undirected.

The number of undirected edges is determined by the Fano predicate.
For $n = 3$ factors there are $C_2 = 2$ parenthesisations and a
single reassociation edge; this edge is undirected iff the triple
is associative. The first nontrivial mixed picture appears at $n = 4$
factors ($C_3 = 5$ parenthesisations, 5 edges), where some edges are
directed and others are not, depending on the Fano status of the
constituent triples.

\subsection{Skew Monoidal Coherence}\label{sec:coherence}

We now prove that the Fano-selective associator structure carries a
full skew monoidal coherence in the sense of Lack--Street
\cite{lackstreet2012}. The five axioms reduce to two verifiable
conditions: the pentagon (3-cocycle) and the triangle (unitor)
identities.

\paragraph{The reassociation factor.}
Define $\varphi(a,b,c) := \alpha(a,b,c) \cdot \beta(a,b,c)$, where
$\alpha$ and $\beta$ are the left and right parenthesisation signs.
Since $\alpha, \beta \in \{+1, -1\}$, we have
$\varphi \in \{+1, -1\}$, with $\varphi = 1$ iff the triple is
associative (Fano).

\paragraph{Unitors.}
The sign function satisfies $\sigma(0, a) = \sigma(a, 0) = 1$ for
all $a \in \{0, \ldots, 7\}$, so the identity element $e_0$ is both
a left and right unit. The left and right unitors $\lambda_a$ and
$\rho_a$ are identities. The four triangle axioms involving $\lambda$,
$\rho$, and $\alpha$ are therefore trivially satisfied.

\paragraph{The pentagon.}
The pentagon axiom for a skew monoidal category requires, for all
$(a,b,c,d)$:
\begin{equation}\label{eq:pentagon}
  \varphi(a \oplus b, c, d) \cdot \varphi(a, b, c \oplus d)
  \;=\;
  \varphi(a, b, c) \cdot \varphi(a, b \oplus c, d) \cdot \varphi(b, c, d)
\end{equation}
This is exactly the \emph{3-cocycle condition} on $\ZZ_2^3$ for the
function $\varphi$. Albuquerque and Majid \cite{albuquerque1999}
showed that the Cayley--Dickson sign function determines a normalised
3-cocycle on $\ZZ_2^k$ for all $k$; equation~\eqref{eq:pentagon} is a
direct consequence.

\paragraph{The ambient category.}
We work in the category $\mathcal{V}$ of $\ZZ_2^3$-graded
one-dimensional real vector spaces: objects are the 8 spaces
$V_g = \RR \cdot e_g$ ($g \in \ZZ_2^3$), tensor is induced by
grade addition ($V_g \otimes V_h = V_{g \oplus h}$) with the
Cayley--Dickson sign, unit is $V_0$, and the associator on
homogeneous elements is multiplication by $\varphi(g,h,k)$.

\begin{theorem}[Cocycle Coherence]\label{thm:coherence}
The category $\mathcal{V}$ with tensor $\oplus$ (XOR), unit $V_0$,
and associator factor $\varphi$ satisfies all five monoidal coherence
axioms. In particular, $\varphi$ is a normalised 3-cocycle on
$\ZZ_2^3$, and the associator is an isomorphism on every triple.
\end{theorem}

\begin{proof}
The unitor identities $\sigma(0,a) = \sigma(a,0) = 1$ are verified
for all $64$ pairs (theorem \texttt{left\_unitor\_identity},
\texttt{right\_unitor\_identity} in \texttt{SounioSkewCategory.lean}).
The pentagon~\eqref{eq:pentagon} is verified for all
$8^4 = 4096$ quadruples (theorem \texttt{pentagon\_coherence}).
The four triangle axioms follow from $\lambda = \rho = \mathrm{id}$.
Since $\varphi \in \{+1,-1\}$ and $(-1)^{-1} = -1$, every associator
is invertible; $\mathcal{V}$ is a (non-symmetric) monoidal category
with a $\ZZ_2^3$ cocycle twist, as in Albuquerque--Majid
\cite{albuquerque1999}.
\end{proof}

\subsection{The Rewrite Category}\label{sec:rewrite-cat}

\begin{remark}\label{rem:semantic}
The algebraic associator is everywhere invertible, but this does
\emph{not} mean every reassociation is a valid compiler
transformation. We now define the category in which the selective
structure lives.
\end{remark}

\begin{definition}[Rewrite category]\label{def:rewrite-cat}
Let $\mathcal{R}$ be the category whose objects are fully
parenthesised expressions over octonion basis elements, and whose
morphisms are \emph{semantics-preserving rewrites}: reassociation
steps $(e_i \cdot e_j) \cdot e_l \to e_i \cdot (e_j \cdot e_l)$
that do not change the computed value. A morphism exists from the
left parenthesisation to the right iff $\varphi(i,j,l) = +1$.
\end{definition}

In $\mathcal{R}$, the 175 associative triples have morphisms (and
their inverses, since the identity map is its own inverse); the 168
sign-changing triples have \emph{no morphism at all}---not even a
directed one. This is stronger than a skew monoidal category (which
would have a directed morphism without an inverse); it is a
\emph{partial} monoidal structure where some reassociations are
simply absent.

The \texttt{NonAssoc} effect in the compiler tracks exactly this
distinction: the effect boundary marks the triples for which
$\mathcal{R}$ has no morphism, preventing the e-graph from
introducing rewrites that would change the program's semantics.

% ================================================================
\section{Implementation in Sounio}\label{sec:impl}
% ================================================================

Sounio is a self-hosted programming language for epistemic and
scientific computing. Its compiler pipeline includes an algebraic
effect system and an e-graph--based optimiser. We implement
Fano-selective invertibility at two levels: the type checker and the
optimiser.

\subsection{The \texttt{NonAssoc} Algebraic Effect}

We add \texttt{NonAssoc} as effect ID~14 in the effect system. Any
function that multiplies non-associative hypercomplex values must
declare this effect:

\begin{lstlisting}[language=Sounio]
fn octonion_product(a: Hyper<Octonion, f64>,
                    b: Hyper<Octonion, f64>)
    -> Hyper<Octonion, f64> with NonAssoc {
  a * b
}
\end{lstlisting}

The type checker enforces the effect requirement:
\begin{lstlisting}[language=Sounio]
// In check_hyper_binary (self-hosted/check/hyper.sio):
if !is_associative(result_alg) {
    if op == BinaryOp::OpMul {
        // E035: requires NonAssoc effect
    }
}
\end{lstlisting}

The effect triggers for all Cayley--Dickson algebras at tower level
$k \geq 3$ (octonions, sedenions, and beyond).

\subsection{The Fano Optimisation Predicate}

The compiler's reassociation optimiser uses the Fano predicate to
decide safety:

\begin{lstlisting}[language=Sounio]
// In ir_can_reassociate_triple (self-hosted/ir/algebra.sio):
fn ir_can_reassociate_triple(bits: i64, i: i64,
    j: i64, k: i64) -> bool with Mut, Panic, Div {
  if bits < 3 { return true }
  if i == 0 || j == 0 || k == 0 { return true }
  if i == j || j == k || i == k { return true }
  let alpha = ir_cd_sigma(i, j, bits)
            * ir_cd_sigma(i ^ j, k, bits)
  let beta = ir_cd_sigma(j, k, bits)
           * ir_cd_sigma(i, j ^ k, bits)
  alpha == beta  // Fano line iff signs agree
}
\end{lstlisting}

The e-graph saturation engine checks this predicate before applying
the associativity rewrite rule. When \texttt{algebra\_kind $\geq 3$}
(non-associative), the guard blocks the rewrite unless
\texttt{ir\_can\_reassociate\_triple} returns true.

This makes $\PG(2,2)$ an operational component of the compiler: the
7 Fano lines determine the 42 distinct-index safe reassociations;
together with the repeated-index cases (guaranteed by alternativity),
these yield the 175 associative triples out of 343.

\begin{remark}[Hard block vs.\ negation injection]
An equality saturation engine could, in principle, apply the
non-Fano reassociation and inject a compensating unary negation
node into the e-class. We deliberately choose a \emph{hard block}
instead: in the category of semantic rewrites (Remark~\ref{rem:semantic}),
a sign-changing transformation is not an identity but an effect
boundary. Injecting arithmetic to compensate would silently
introduce an operation that the programmer did not write and that
the \texttt{NonAssoc} effect is designed to make visible. The hard
block preserves the invariant that every rewrite in the e-graph is
semantics-preserving \emph{without} auxiliary compensation.
\end{remark}

\subsection{Artin's Theorem as a Compiler Optimisation}

For alternative algebras ($k = 3$, octonions), Artin's theorem
guarantees that any subalgebra generated by two elements is
associative. The compiler implements this:

\begin{lstlisting}[language=Sounio]
fn ir_artin_safe(alg: i64,
    n_distinct_vars: i64) -> bool {
  if alg <= 2 { return true }
  if alg == 3 && n_distinct_vars <= 2
    { return true }
  return false
}
\end{lstlisting}

This is a coarser but faster check: if an expression involves at most
two distinct variables, the compiler can reassociate freely even in the
octonions. Note that \texttt{n\_distinct\_vars $\leq$ 2} captures
exactly the 133 repeated-index triples from the partition table in
\S\ref{sec:background}---the triples that are associative by
alternativity (Artin's theorem), independent of any Fano-line
geometry. The Fano predicate provides finer-grained control for
expressions with three or more distinct variables, recovering the
remaining 42 associative triples from the 7 Fano lines.

% ================================================================
\section{Formal Verification in Lean~4}\label{sec:formal}
% ================================================================

All results are formalised in Lean~4 \cite{moura2021lean4} across
four files totalling approximately 1,250 lines.

\subsection{Architecture}

\begin{center}
\begin{tabular}{lrl}
\toprule
File & Theorems & Build time \\
\midrule
\texttt{SounioCayleyDickson.lean} & 12 & 666\,ms \\
\texttt{SounioSkewCategory.lean} & 19 & 672\,ms \\
\texttt{SounioBidirectionalBridge.lean}$^*$ & 18 & 420\,ms \\
\texttt{SounioEffects.lean} (modified) & 46 & 621\,ms \\
\midrule
Total & 95 & 2.4\,s \\
\bottomrule
\end{tabular}
\end{center}

No \texttt{sorry}. No Mathlib dependency. All counting theorems are
proved by \texttt{native\_decide}, exploiting the decidability of
Boolean predicates over finite domains.
($^*$\texttt{SounioBidirectionalBridge.lean} proves that the
$\alpha/\beta$ sign decomposition, when labelled as
``synthesis/checking'', reproduces the same tripartition. This is a
definitional relabelling, not an independent result; we include it for
completeness and cross-validation.)

\subsection{The Sign Function in Lean~4}

The Lean formalisation mirrors the Sounio implementation:

\begin{lstlisting}[language=Lean4]
def cdSigma (a b bits : Nat) : Int :=
  if a = 0 || b = 0 then 1
  else if bits <= 1 then -1
  else
    let half := Nat.shiftLeft 1 (bits - 1)
    let aHi := a >= half
    let bHi := b >= half
    let aLo := a &&& (half - 1)
    let bLo := b &&& (half - 1)
    if !aHi && !bHi then cdSigma aLo bLo (bits-1)
    else if !aHi && bHi then cdSigma bLo aLo (bits-1)
    else if aHi && !bHi then
      if bLo = 0 then cdSigma aLo bLo (bits-1)
      else -(cdSigma aLo bLo (bits - 1))
    else
      if bLo = 0 then -(cdSigma bLo aLo (bits-1))
      else cdSigma bLo aLo (bits - 1)
termination_by bits
\end{lstlisting}

\subsection{Cross-Validation}

The Lean formalisation and the Sounio implementation are
cross-validated by three runtime verification programs
(\texttt{examples/categorical\_bridge\_*.sio}) that independently
compute all 343 triples and verify the same counts. The theorem
\texttt{compiler\_predicate\_match} in Lean formally states that the
Fano predicate agrees with the associator-coefficient-zero predicate for all
343 triples.

% ================================================================
\section{Structural Analogy to the Two-State Vector Formalism}
\label{sec:tsvf}
% ================================================================

We briefly note a structural analogy with the two-state vector
formalism (TSVF) of Aharonov, Bergmann, and Lebowitz
\cite{abl1964,aharonov2008tsvf}.
In the TSVF, a quantum system is described by a forward-evolving state
$|\psi\rangle$ and a backward-evolving state $\langle\phi|$; their
interference determines observables. The Fano-selective structure
exhibits the same combinatorics: $\alpha$ (left parenthesisation) and
$\beta$ (right parenthesisation) play the roles of forward and backward
signs, the associator coefficient $\alpha - \beta$ is the interference, the
dagger $(i,j,l) \mapsto (l,j,i)$ exchanges forward and backward with
perfect $84/84$ symmetry, and Fano triples correspond to
``time-symmetric'' observables where both directions agree.

This analogy is \emph{structural}, not physical. We do not claim that
the compiler performs retrocausal computation. Whether the shared
tripartition structure reflects a deeper connection through compact
closed categories with daggers \cite{abramsky2004} is an open question.

% ================================================================
\section{Related Work}\label{sec:related}
% ================================================================

\paragraph{Skew monoidal categories.}
Szlach\'{a}nyi \cite{szlachanyi2012} and Lack--Street
\cite{lackstreet2012} introduced skew monoidal categories to study
bialgebroids. Bourke--Lack \cite{bourke2018} connected them to
skew multicategories. Buckley et al.\ \cite{buckley2014} studied the Catalan simplicial set
as a classifying structure for skew monoidal categories.
Uustalu, Veltri, and Zeilberger \cite{uvz2021} studied partially
normal skew monoidal categories, where selected structural laws become
invertible---the closest precedent to our work. However, their
invertibility conditions are uniform (e.g., all left unitors are
invertible), not indexed by a geometric predicate on basis triples.

\paragraph{Dagger categories and PL.}
Sabry et al.\ \cite{sabry2024quantumpi} implement daggers as a
language feature in $\Pi$/Quantum$\Pi$. Karvonen \cite{karvonen2019}
developed dagger category theory systematically. Heunen--Karvonen
\cite{heunen2016} studied dagger monads. None of these works connect
daggers to non-associative algebra or the Fano plane.

\paragraph{Octonions as quasialgebra.}
Albuquerque--Majid \cite{albuquerque1999} showed the octonions are an
algebra object in a monoidal category with a non-trivial $\ZZ_2^3$
3-cocycle associator. This resolves non-associativity by moving to a
twisted monoidal category. Our approach preserves non-associativity as
a first-class feature, selectively tracked by the effect system.

\paragraph{Non-associative Lambek calculus.}
Lambek \cite{lambek1961} introduced the non-associative calculus NL
with binary tree sequents. Buszkowski established polynomial-time
decidability. A recent Agda formalisation \cite{agdanl2025} proves cut
admissibility and interpolation. NL is a substructural logic, not a
type system with algebraic effects; we share the binary tree structure
but add the Fano selection and effect tracking.

\paragraph{Bidirectional type checking.}
Dunfield--Krishnaswami \cite{dunfield2021} survey bidirectional typing.
Mihejevs--Hedges \cite{mihejevs2025} identify a three-way
polarity/chirality/mode correspondence. The synthesis/checking modes
provide a natural two-directional structure analogous to our
$\alpha/\beta$ decomposition, but we make no deep claim about this
connection; it is a labelling, not a theorem.

\paragraph{Formal verification of algebra.}
Cowles--Gamboa \cite{cowles2017} verified Cayley--Dickson properties
in ACL2. Our work provides the first Lean~4 formalisation of the
168 theorem, the binary norm, and the dagger decomposition.

% ================================================================
\section{Conclusion}\label{sec:conclusion}
% ================================================================

We have introduced Fano-selective invertibility, a categorical
structure in which the geometry of the Fano plane determines which
associator morphisms are invertible. This structure:

\begin{itemize}
\item Generalises both monoidal and skew monoidal categories.
\item Is instantiated by the Cayley--Dickson octonion algebra with
  the exact partition $175 + 84 + 84 = 343$.
\item Is implemented as the \texttt{NonAssoc} algebraic effect and
  the Fano-based reassociation predicate in a working self-hosted compiler.
\item Is formalised in Lean~4 with 95 machine-checked theorems.
\end{itemize}

\paragraph{Open questions.}
\begin{enumerate}
\item Does Fano-selective invertibility extend to a full
  \emph{Fano-selective skew closed category}, with non-invertible
  evaluation and currying maps for non-Fano triples?
\item For higher Cayley--Dickson algebras ($k \geq 4$), what is the
  structure of the Fano-selective partition? The ratio of non-Fano to
  total triples approaches $1/2$ as $k \to \infty$ (Corollary~C of
  \cite{agourakis2026}).
\item Is the structural analogy to TSVF an accident of small numbers,
  or does it reflect a deeper connection through compact closed
  categories with daggers?
\end{enumerate}

\paragraph{Artifact.}
We freeze the self-hosted compiler artifact at release tag
\texttt{self-host-native-v1} (commit \texttt{1aef4e65}). At this
checkpoint the native compiler reaches a bit-identical fixed point
(\texttt{gen2 == gen3}, MD5 \texttt{7448f33f5a5024fbada0c70a75e616ce}),
passes $82/82$ native runtime proofs, $16/16$ native typecheck proofs,
and $4/4$ ARM64 compile proofs, with no Rust JIT required in the
validation path. The Lean~4 formalisation, Sounio compiler, and runtime
verification examples are available at
\url{https://github.com/dagourakis/sounio}.

% ================================================================
% References
% ================================================================

\bibliographystyle{ACM-Reference-Format}
\begin{thebibliography}{99}

\bibitem{abl1964}
Y.~Aharonov, P.G.~Bergmann, and J.L.~Lebowitz.
\newblock Time symmetry in the quantum process of measurement.
\newblock {\em Physical Review}, 134:B1410--B1416, 1964.

\bibitem{aharonov2008tsvf}
Y.~Aharonov and L.~Vaidman.
\newblock The two-state vector formalism: An updated review.
\newblock In {\em Time in Quantum Mechanics}, volume 734 of {\em Lecture Notes
  in Physics}, pages 399--447. Springer, 2008.

\bibitem{abramsky2004}
S.~Abramsky and B.~Coecke.
\newblock A categorical semantics of quantum protocols.
\newblock In {\em LICS 2004}, pages 415--425. IEEE, 2004.

\bibitem{agdanl2025}
N.~Veltri and C.-S.~Wan.
\newblock An {A}gda formalisation of nonassociative {L}ambek calculus and its metatheory.
\newblock In {\em TABLEAUX 2025}, volume 15980 of {\em LNAI}, pages 433--452. Springer, 2025.

\bibitem{agourakis2026}
D.~Agourakis.
\newblock The 168 theorem: Nonzero associators in {C}ayley--{D}ickson algebras.
\newblock {\em Advances in Applied Clifford Algebras} (submitted), 2026.

\bibitem{albuquerque1999}
H.~Albuquerque and S.~Majid.
\newblock Quasialgebra structure of the octonions.
\newblock {\em Journal of Algebra}, 220(1):188--224, 1999.

\bibitem{baez2002octonions}
J.C.~Baez.
\newblock The octonions.
\newblock {\em Bulletin of the American Mathematical Society (N.S.)},
  39(2):145--205, 2002.

\bibitem{bourke2018}
J.~Bourke and S.~Lack.
\newblock Skew monoidal categories and skew multicategories.
\newblock {\em Journal of Algebra}, 506:237--266, 2018.

\bibitem{buckley2014}
M.~Buckley, R.~Garner, S.~Lack, and R.~Street.
\newblock The {C}atalan simplicial set.
\newblock {\em Math.\ Proc.\ Cambridge Philos.\ Soc.}, 158(2):211--222, 2015.

\bibitem{cowles2017}
J.~Cowles and R.~Gamboa.
\newblock Cayley--{D}ickson and {C}lifford algebras in {ACL2}.
\newblock In {\em ACL2 Workshop}, 2017.

\bibitem{dickson1919}
L.E.~Dickson.
\newblock On quaternions and their generalization and the history of the
  eight-square theorem.
\newblock {\em Annals of Mathematics}, 20(3):155--171, 1919.

\bibitem{dunfield2021}
J.~Dunfield and N.~Krishnaswami.
\newblock Bidirectional typing.
\newblock {\em ACM Computing Surveys}, 54(5), Article~98, 1--38, 2021.

\bibitem{heunen2016}
C.~Heunen and M.~Karvonen.
\newblock Monads on dagger categories.
\newblock {\em Theory and Applications of Categories}, 31(35):1016--1043, 2016.

\bibitem{karvonen2019}
M.~Karvonen.
\newblock {\em The Way of the Dagger}.
\newblock PhD thesis, University of Edinburgh, 2019.

\bibitem{lackstreet2012}
S.~Lack and R.~Street.
\newblock Skew monoidales, skew warpings and quantum categories.
\newblock {\em Theory and Applications of Categories}, 26(15):385--402, 2012.

\bibitem{lambek1961}
J.~Lambek.
\newblock On the calculus of syntactic types.
\newblock In R.~Jakobson, editor, {\em Structure of Language and Its
  Mathematical Aspects}, pages 166--178. AMS, 1961.

\bibitem{leijen2014}
D.~Leijen.
\newblock Koka: Programming with row polymorphic effect types.
\newblock In {\em HOPE 2014}, 2014.

\bibitem{lindley2012}
S.~Lindley and J.~Cheney.
\newblock Row-based effect types for database integration.
\newblock In {\em TLDI 2012}, 2012.

\bibitem{mihejevs2025}
V.~Mihejevs and J.~Hedges.
\newblock Canonical bidirectional typechecking.
\newblock arXiv:2512.07511, 2025.

\bibitem{moura2021lean4}
L.~de~Moura and S.~Ullrich.
\newblock The {L}ean 4 theorem prover and programming language.
\newblock In {\em CADE-28}, volume 12699 of {\em LNAI}, pages 625--635, 2021.

\bibitem{plotkin2009}
G.~Plotkin and M.~Pretnar.
\newblock Handlers of algebraic effects.
\newblock In {\em ESOP 2009}, volume 5502 of {\em LNCS}, pages 80--94, 2009.

\bibitem{renzhao2023}
G.~Ren and Y.~Zhao.
\newblock The explicit twisted group algebra structure of the {C}ayley--{D}ickson algebra.
\newblock {\em Advances in Applied Clifford Algebras}, 33, Article~49, 2023.

\bibitem{sabry2024quantumpi}
J.~Carette, C.~Heunen, R.~Kaarsgaard, and A.~Sabry.
\newblock How to bake a quantum pi.
\newblock {\em Proc.\ ACM Program.\ Lang.}, 8(ICFP), Article~236, 1--29, 2024.

\bibitem{selinger2007}
P.~Selinger.
\newblock Dagger compact closed categories and completely positive maps.
\newblock {\em Electronic Notes in Theoretical Computer Science}, 170:139--163, 2007.

\bibitem{szlachanyi2012}
K.~Szlach\'{a}nyi.
\newblock Skew-monoidal categories and bialgebroids.
\newblock {\em Advances in Mathematics}, 231(3--4):1694--1730, 2012.

\bibitem{tamari1962}
D.~Tamari.
\newblock The algebra of bracketings and their enumeration.
\newblock {\em Nieuw Archief voor Wiskunde}, 10:131--146, 1962.

\bibitem{uvz2021}
T.~Uustalu, N.~Veltri, and N.~Zeilberger.
\newblock Proof theory of partially normal skew monoidal categories.
\newblock {\em EPTCS}, 333:230--246, 2021.

\end{thebibliography}

\end{document}
